\section{Conclusion}
\label{sec:conclusion}

Task--nuisance ambiguity has a classification cost when statistically similar alternatives require different decisions. Our local bound makes that condition explicit, and a worked family shows how nuisance confounding reduces the effective sample size near a decision boundary despite constant raw information at the reference parameter. A matching upper bound identifies the learning rate and the cap imposed by the local decision scale. Population audits then account for estimation error, ridge bias, adaptive membership, and support, establishing which information statements are supported by data.

The retained Waterbirds experiments establish a 52.6-point controlled accuracy drop in a frozen predefined group and incremental species-held-out failure-prediction value from nuisance-derived coordinates beyond a specified confidence baseline. Independent evaluation validates a discovered low predictive-information population. Training and acquisition show that information gains and predictive gains can diverge.

The classification result, population audit, and behavioral evaluations have complementary roles. They connect ambiguity to learning difficulty, quantify a declared information target with uncertainty, and test its relationship to observed behavior. This supports evaluating information objectives together with predictive performance and controlled nuisance sensitivity when assessing a classifier or a proposed improvement.
